\chapter{Memoria Descriptiva}

\section{Generalidades del Proyecto}
El presente Expediente Técnico de Ingeniería corresponde al diseño, dimensión, especificación y sustento normativo de las **Instalaciones Eléctricas Industriales** para una nave industrial de producción y almacenamiento de $20\text{ m}$ de ancho por $40\text{ m}$ de largo ($800\text{ m}^2$ de área techada total, con una altura libre de $8\text{ m}$).

El proyecto contempla la acometida trifásica en Baja Tensión ($380\text{V} / 220\text{V}$ a $60\text{ Hz}$), el Tablero General de Distribución (TGD), sub-tableros de distribución de fuerza (TF1) e iluminación (TI1), sistema de iluminación industrial mediante luminarias LED High-Bay de alta eficiencia a $7\text{ m}$ de altura, tomacorrientes industriales trifásicos tipo Stecker y monofásicos Shucko, fuerza motriz para un puente grúa de 5Tn ($10\text{ HP}$), compresor de aire industrial de $15\text{ HP}$ con arranque Estrella-Delta, maquinaria pesada de taller ($25\text{ kW}$), servicios administrativos de oficinas y un banco automático de condensadores de $15\text{ kVAr}$ para el control del factor de potencia.

\section{Ubicación Geográfica y Política}
\begin{itemize}
    \item **Departamento:** Puno
    \item **Provincia:** San Román
    \item **Distrito:** Juliaca
    \item **Dirección:** Zona Industrial de Juliaca S/N
    \item **Altitud:** $3,824\text{ m.s.n.m.}$
\end{itemize}

\section{Alcance de las Obras}
El alcance del presente proyecto comprende la ejecución de las siguientes partidas principales:
\begin{enumerate}
    \item Acometida eléctrica trifásica en cable N2XH $50\text{ mm}^2$ desde la red de concesionaria hasta el medidor y TGD.
    \item Suministro e instalación del Tablero General de Distribución TGD-Nave (24 polos, IP65) y sub-tableros TF1 y TI1.
    \item Sistema de iluminación industrial con 20 campanas LED High-Bay de $150\text{W}$ 5000K IP65 y paneles LED en oficinas.
    \item Red de tomacorrientes industriales Stecker de $32\text{A}$ $380\text{V}$ (3F+PE) y tomas monofásicos de $16\text{A}$ $220\text{V}$.
    \item Circuito de fuerza motriz para Puente Grúa de $10\text{ HP}$ (regimen intermitente) con guardamotor regulable.
    \item Circuito de fuerza motriz para Compresor de $15\text{ HP}$ con temporizador y contactores para arranque Estrella-Delta.
    \item Circuito de fuerza para maquinaria pesada de taller ($25\text{ kW}$) compuesto por tornos, fresadoras y cortadora de plasma.
    \item Compensación automática de energía reactiva mediante banco de condensadores de $15\text{ kVAr}$ en 3 pasos ($3 \times 5\text{ kVAr}$).
    \item Sistema de Puesta a Tierra (SPAT) perimetral compuesto por malla enterrada a $0.80\text{ m}$ con 6 electrodos Copperweld de $5/8'' \times 2.40\text{ m}$ e interconexión en cable de cobre desnudo de $35\text{ mm}^2$.
\end{enumerate}

\section{Condiciones del Suministro Eléctrico}
\begin{table}[H]
\centering
\caption{Parámetros fundamentales del suministro eléctrico industrial}
\begin{tabular}{ll}
\toprule
\textbf{Parámetro} & \textbf{Valor Especificado} \\
\midrule
Tensión Nominal de Red & $380\text{ V}$ (entre fases) / $220\text{ V}$ (fase-neutro) \\
Sistema Eléctrico & Trifásico $3\text{F} + \text{N} + \text{PE}$ \\
Frecuencia Nominal & $60\text{ Hz}$ \\
Esquema de Conexión a Tierra & TN-S (Neutro y PE separados desde TGD) \\
Potencia Instalada Total (P.I.) & $55.20\text{ kW}$ \\
Máxima Demanda Adoptada (M.D.) & $39.35\text{ kW}$ \\
Factor de Simultaneidad Aplicado & $0.85$ (según CNE-Utilización) \\
Reserva de Crecimiento & $20\%$ \\
Corriente Total de Trabajo ($I_b$) & $66.44\text{ A}$ \\
Factor de Potencia Objetivo & $\cos\phi \ge 0.96$ \\
\bottomrule
\end{tabular}
\end{table}

\section{Reglamento y Normas de Aplicación}
El diseño y los cálculos han sido desarrollados en estricto cumplimiento de las siguientes normas técnicas:
\begin{itemize}
    \item **Código Nacional de Electricidad - Utilización (CNE-U 2006)**: Reglas 050-200, Regla 050-202, Tablas 4 y 5.
    \item **Reglamento Nacional de Edificaciones (RNE)**: Norma Técnica EM.010 - Instalaciones Eléctricas Interiores.
    \item **Normas DGE / MEM**: Simbología unificada para planos eléctricos (Resolución Ministerial N° 091-2002-EM/VME).
    \item **Reglamento Nacional de Metrados**: Norma Técnica OE.5 - Instalaciones Eléctricas y Mecánicas.
    \item **IEC 60364**: Low-voltage electrical installations.
    \item **IEC 60947-2**: Low-voltage switchgear and controlgear - Circuit-breakers.
    \item **IEEE Std 399 (Brown Book)**: Industrial and Commercial Power Systems Analysis (Arranque de motores y caída de tensión).
\end{itemize}
